\begin{onehalfspace}

I risultati delimitano con chiarezza ciò che la demo e la campagna
dimostrano: l'esecuzione ideale è verificata anche per $N=21/35$, mentre
il confronto controllato fra canali di rumore è sostenibile su $N=15$.
Le estensioni seguenti mirano a colmare questa distanza senza
estrapolare i risultati oltre il loro perimetro.

% -----------------------------------------------
\section{Scalabilità a Istanze di Maggiore Dimensione}
\label{sec:scalabilita_futura}

La priorità è ridurre il costo dell'aritmetica modulare. La costruzione
Beauregard implementata nella tesi è stata validata idealmente mediante
truth table e distribuzione QPE, ma per $N=21$, $a=2$ e otto controlli il
circuito esecutivo compilato nella base \texttt{rz/sx/x/cx} contiene
21\,036 CX e ha profondità 23\,081. Questi sono conteggi della
compilazione congiunta, non somme di blocchi né stime su un gate set più
astratto.

Le direzioni immediate sono l'approximate QFT con soglia d'angolo
esplicita, addizionatori alternativi, sintesi mirata al gate set del
backend e riduzione delle operazioni controllate prima della
traspilazione. Ogni variante dovrebbe superare gli stessi controlli già
usati qui: verità funzionale della moltiplicazione, ancille pulite,
coerenza di fase e distanza della QPE ideale dalla legge teorica. Il
numero di gate, da solo, non è una prova di correttezza.

Anche la scelta delle istanze deve essere governata dall'ordine
moltiplicativo $r$, non soltanto dal modulo $N$. L'istanza $N=35$, $a=6$
ha ordine 2 e non costituisce una progressione monotona rispetto a
$N=21$, $a=2$, che ha ordine 6. Una campagna di scalabilità dovrebbe
quindi stratificare sia la larghezza sia l'ordine e dichiarare a priori
quale delle due variabili intende studiare.

% -----------------------------------------------
\section{Ottimizzazione della Trasformata di Fourier Consapevole della Topologia}
\label{sec:qft_topologia}

La direzione precedente indica l'approximate QFT fra le varianti da esplorare.
Una campagna esplorativa condotta a chiusura del lavoro ne ha misurato il margine, e il
risultato merita di essere riportato perché delimita l'intera famiglia di ottimizzazioni
di questo tipo.

\paragraph{Controllo su $N=15$ ed eliminazione degli SWAP.}
Nel pilota su $N=15$, la configurazione $k_{\mathrm{QPE}}=1$, scelta sui
batch di selezione, ottiene sull'holdout $0.6279$ contro $0.4792$ della
QFT piena: differenza $+0.1487$, IC Newcombe 95\%
$[+0.1273;+0.1699]$, con 279 ECR invece di 362. Il beneficio riguarda
un caso favorevole e degenere, con fasi esatte in due bit, e non dimostra
lo stesso effetto per ordini più lunghi. L'eliminazione degli SWAP di
inversione dei bit produce invece conteggi compilati e risultati
identici: il transpiler ne assorbe già la permutazione nel layout.
Griglia, controllo della variante corretta e risultati completi sono
nella Sezione~\ref{sec:app_aqft} dell'Appendice~\ref{app:fase2}.

Nell'aritmetica di Beauregard, gran parte del costo viene dalle QFT interne agli addizionatori modulari. Per $N=21$ si contano 11\,436 porte \texttt{cp} e 3\,406 \texttt{cx}; la QFT finale richiede invece solo 28 rotazioni, su circa 47\,000 porte complessive. Eliminare le rotazioni più piccole negli addizionatori porta le porte a due qubit da 49\,661 a 23\,178. Agire sulla sola trasformata finale non cambia l'ordine di grandezza. Il primo controllo identifica quindi la parte del circuito su cui conviene intervenire.

Ridurre le porte comporta una perdita di precisione. Su $N=21$, il successo ideale scende da $0{,}4667$ a $0{,}4013$, cioè di circa sei punti e mezzo. In presenza di rumore, la variante troncata può recuperare questa perdita solo se il risparmio di porte evita abbastanza errori. Il bilancio dipende quindi dalla qualità delle porte e va verificato.

Un modello semplificato combina il successo ideale con quello di una distribuzione uniforme. In questo modello, il vantaggio massimo stimato della variante troncata è circa \textbf{tre punti percentuali}. Richiederebbe errori a due qubit circa \textbf{duecento volte inferiori} a quelli dei dati di calibrazione considerati. È una stima del modello, non un vantaggio misurato né un limite universale. Nei test rumorosi non emerge un beneficio statisticamente significativo del troncamento. Al rumore nominale, le configurazioni provate restano vicine al livello uniforme; nel punto in cui compare un segnale, è la QFT \emph{piena} ad avere il risultato migliore. Il circuito di $N=21$ dura circa centocinquanta volte il tempo di coerenza tipico. Questa profondità limita la possibilità di distinguere le varianti nel regime considerato.

Per capire se il limite dipenda dal costo di Shor, si ripete il confronto sulla sola \textbf{stima di fase}. Con un unitario di fase, il circuito richiede decine di porte a due qubit, anziché decine di migliaia. I nove confronti usano tre dimensioni di registro e fasi sia esattamente rappresentabili sia non rappresentabili in $n$ bit. La variante troncata ha un risultato migliore in \textbf{tutti e nove}; in sei, l'intervallo di confidenza è interamente positivo. Su dieci qubit il vantaggio raggiunge $+0{,}126$, con IC Newcombe 95\% $[+0{,}098;\ +0{,}153]$.

Nei test sulla stima di fase emergono tre regolarità. Il grado migliore di troncamento resta \textbf{due o tre} alle taglie provate, mentre la QFT piena usa $n-1$. Il risparmio cresce così da venticinque porte a due qubit su sei qubit a settantotto su dieci. Cresce anche il vantaggio, da circa due a tredici punti percentuali. Infine, le fasi esattamente rappresentabili beneficiano più delle altre. Questo aiuta a interpretare la differenza fra le istanze di Shor: per $N=15$, l'ordine $r=4$ dà fasi $0$, $1/4$, $1/2$, $3/4$, esatte in due bit. Per $N=21$, l'ordine $r=6$ dà invece fasi $j/6$, per le quali le rotazioni fini conservano informazione utile.

La prima di queste regolarità --- quella che a prima vista sorprende di più, perché un
registro più lungo sembrerebbe richiedere più rotazioni fini --- ha una spiegazione
teorica indipendente, pubblicata mentre questo lavoro era in corso. Akoramurthy e
Surendiran derivano per la stima di fase con trasformata troncata un limite in forma
chiusa sulla distanza in variazione totale fra la distribuzione esatta e quella
troncata~\cite{akoramurthy2026}:
%
\begin{equation}
\mathrm{TV}\!\left(P_\phi,\, P_\phi^{d}\right) \;\le\; \frac{\pi\,(m-d)}{2^{d}},
\label{eq:bound_troncamento}
\end{equation}
%
dove $m$ è la dimensione del registro e $d$ il grado di troncamento. Il numeratore cresce
linearmente in $m$, il denominatore esponenzialmente in $d$: per mantenere
fisso il limite al crescere del registro occorre far crescere $d$ circa
logaritmicamente. Il grado piccolo osservato alle tre taglie provate
non dimostra dunque che un grado costante basti a qualunque dimensione. Gli autori
riportano inoltre una riduzione del $31{,}3$--$43{,}7\%$ nel conteggio di porte e un
passaggio di complessità da $O(m^2)$ a $O(m \log m)$ per $d = O(\log m)$. Il loro impianto
sperimentale si ferma però alla stima di fase e non affronta la fattorizzazione: la
delimitazione riportata sopra --- che su Shor la stessa leva vale al massimo tre punti
percentuali, e solo con porte a due qubit molto più accurate delle attuali --- resta
propria di questo lavoro, così come la selezione del grado su dati disgiunti da quelli di
verifica.

Conteggi di porte e calibrazione permettono di proporre gradi candidati
prima dell'esecuzione. In questa campagna, però, il grado viene scelto
empiricamente sui batch di selezione e verificato su holdout disgiunti:
non è stata validata una scelta ottima dai soli conteggi. Il contributo
è quindi un criterio di progetto e un protocollo di verifica per la
trasformata, da controllare quando cambia l'algoritmo o il dispositivo.

Un'estensione consiste nel combinare tre riduzioni. La prima, misurata qui, approssima la QFT per usare meno \emph{porte}. La seconda approssima l'aritmetica modulare per usare meno \emph{qubit di lavoro}~\cite{chevignard2025} e contribuisce alle recenti stime di risorse~\cite{gidney2025}. La Tabella~\ref{tab:traiettoria_risorse} ne riporta l'evoluzione e le assunzioni. La terza divide la stima di fase in blocchi indipendenti per ridurre i \emph{qubit di conteggio}~\cite{shukla2026}.

I tre risparmi non si sommano automaticamente. Ridurre i qubit di lavoro può richiedere più porte: per RSA-2048, una proposta passa da 6144 a 1730 qubit logici, ma porta i Toffoli da $2^{31{,}3}$ a $2^{40{,}9}$~\cite{chevignard2025}. Ottimizzazioni aritmetiche successive recuperano parte di questo costo~\cite{gidney2025}. Combinare la riduzione dei qubit con il troncamento della QFT richiede quindi di misurare il bilancio finale delle porte.

La riduzione del registro di conteggio agisce invece sulla stima di fase e lascia invariato il registro di lavoro~\cite{shukla2026}. La proposta passa da $2n+1$ a tre o quattro qubit di conteggio. Blocchi indipendenti stimano più bit di fase, poi un procedimento dal bit meno significativo al più significativo risolve le ambiguità~\cite{shuklaqpe2026}. Per verificare se questa riduzione sia compatibile con le altre, occorre un esperimento che le combini e misuri anche la profondità. Le singole proposte sono quantificate; il loro beneficio congiunto resta da dimostrare.

La QFT compare anche nella stima di fase, nell'amplitude estimation, nel conteggio quantistico, nella risoluzione di sistemi lineari e nel problema del sottogruppo nascosto, di cui Shor è un caso particolare. Si usa inoltre nelle simulazioni che alternano rappresentazioni in posizione e in momento. Questi algoritmi impiegano le stesse rotazioni controllate di piccolo angolo. Il protocollo seguito qui può suggerire verifiche analoghe: proporre gradi di troncamento dai conteggi, sceglierli su dati di selezione e valutarli su dati indipendenti. Queste verifiche restano da eseguire. Il contributo è quindi un metodo di valutazione, con risultati limitati ai circuiti provati.

Il pilota incompleto su $N=21$, i tre livelli di sensibilità al rumore,
la distinzione fra configurazione scelta e variante troncata, e i nove
contrasti QPE sono documentati nelle Tabelle~\ref{tab:aqft_n21_sensibilita}
e~\ref{tab:aqft_qpe_completa}. Sono riportati anche gli esiti negativi
e i contrasti i cui intervalli includono lo zero.

La verifica usa un layout per dimensione, una sola calibrazione e tre fasi bersaglio. Le varianti condividono batch e semi, quindi il confronto interno usa gli stessi campioni. Il risultato non dimostra però che il beneficio si mantenga su altri dispositivi o topologie.

Il precedente più vicino è la compilazione \emph{noise-adaptive}, che ordina o vincola il
collocamento del circuito usando i dati di calibrazione~\cite{murali2019, tannu2019}. La
differenza rispetto a quella linea, e il motivo per cui resta spazio, è che tali metodi
ottimizzano la fedeltà del circuito, mentre per un algoritmo come Shor conta soltanto che
il post-processing restituisca i fattori. I due obiettivi non coincidono, e la campagna
sulla compilazione riportata in Appendice~\ref{app:fase2} misura quanto non coincidano.

% -----------------------------------------------
\section{Validazione su Hardware Quantistico Reale}
\label{sec:hw_reale}

Va anzitutto dichiarato che l'assenza di esecuzione su hardware è una
\textbf{scelta deliberata e concordata}, non una questione rimasta
irrisolta: pur essendo divenuto disponibile in corso d'opera l'accesso a
un dispositivo, si è convenuto di non impiegarlo, poiché le domande
sperimentali del lavoro erano già state chiuse sui modelli di rumore e
sulle snapshot di calibrazione. L'intera validazione riportata in questa
tesi appartiene quindi a quel perimetro, e ogni confronto con una
macchina reale resta un lavoro successivo, di cui questa sezione descrive
il protocollo.

Il modello rumoroso della tesi è uniforme e illustrativo: non include
coupling map, routing, crosstalk, leakage o deriva della calibrazione. Il
passo successivo è ripetere il protocollo $N=15$ su una QPU, registrando
backend, timestamp, proprietà dei qubit, layout iniziale e finale,
conteggi post-routing e job identifier. Il confronto deve usare lo stesso
budget di shot e una politica di seed dichiarata, senza presentare il
preset uniforme come replica della macchina.

La demo è già predisposta per questo confronto metodologico: distingue
rumore disattivato, scenario moderato, scenario di stress e canali
singoli. Un'estensione hardware dovrebbe affiancare una quinta modalità
``calibrazione acquisita'', costruita dallo snapshot reale e marcata con
la sua data di validità. In questo modo diventerebbe misurabile il
\emph{simulation gap} fra il modello controllato e il dispositivo.

Per $N=21/35$ l'hardware non elimina il problema della profondità: evita
il costo esponenziale dello statevector classico, ma non le migliaia di
operazioni rumorose e i vincoli di connettività. Prima di eseguire queste
istanze servono quindi una sintesi più corta e un budget d'errore basato
sul circuito effettivamente instradato.

% -----------------------------------------------
\section{Generalizzazione ad Altri Algoritmi Quantistici}
\label{sec:altri_algoritmi}

TOP-K è utile soltanto quando più candidati possono essere verificati
classicamente a costo contenuto. Questa proprietà vale per Shor, perché
ogni candidato al periodo produce o meno divisori verificabili, e può
valere per problemi di ricerca con un oracolo di verifica. Non segue
automaticamente che la stessa strategia migliori VQE~\cite{peruzzo2014} o
QAOA~\cite{farhi2014}, dove
l'output centrale è una stima di valore atteso e non una lista di
soluzioni equivalenti.

Una generalizzazione corretta dovrebbe quindi definire prima:
\begin{enumerate}
    \item una funzione classica di validazione del candidato;
    \item il costo aggiuntivo delle $K$ verifiche;
    \item una baseline TOP-1 appaiata sugli stessi campioni;
    \item un criterio che impedisca a $K$ di rendere il successo quasi
    certo per puro effetto combinatorio.
\end{enumerate}
Grover su problemi con verifica economica è un primo banco di prova;
per gli algoritmi variazionali è invece più naturale studiare
allocazione adattiva degli shot o mitigazione del bias dell'osservabile.

% -----------------------------------------------
\section{Integrazione con la Quantum Error Correction}
\label{sec:qec_integrazione}

La campagna di Shor logico usa $p_L$ come probabilità fenomenologica di
Pauli non-identità per gate compilato incluso nel proxy. Non identifica questo parametro con
il logical error rate di un memory experiment surface-code: le due
grandezze hanno circuiti, osservabili e granularità differenti. Il passo
successivo scientificamente significativo è una mappatura esplicita fra
un set di operazioni logiche fault-tolerant e i canali residui stimati per
ciascuna operazione.

Solo dopo tale mappatura sarebbe lecito combinare surface code e Shor in
una previsione end-to-end. Servirebbero almeno: conteggio delle operazioni
logiche per tipo, costo di magic-state distillation per le non-Clifford,
decoder e latenza di feed-forward, errori correlati e intervalli di
confidenza propagati. Il modello fenomenologico della tesi resta utile
come analisi di sensibilità, non come stima di risorse hardware.

Per la decodifica appresa, le direzioni più promettenti restano modelli a
grafo o convoluzionali che rispettino la località del reticolo, uso della
storia temporale dei cicli e soft information analogica della misura.
Queste estensioni aggiungono informazione o struttura; non si limitano a
ingrandire una rete densa sullo stesso vettore di sindrome.

% -----------------------------------------------
\section{Evoluzione del Classificatore ML}
\label{sec:evoluzione_ml}

La campagna aggiornata mostra che TOP-1 è addestrabile ma operativo solo
come predittore della riuscita del candidato dominante; TOP-16 produce
invece $2\,000/2\,000$ positivi in entrambi gli scenari. Inoltre M2 è
peggiore della propria ablazione TOP-4 nei due confronti appaiati. Non è
quindi giustificato migliorare soltanto l'architettura mantenendo la
stessa etichetta.

Un nuovo problema supervisionato dovrebbe prevedere un obiettivo legato
al costo, ad esempio il valore atteso degli shot risparmiati dopo aver
considerato il costo di uno scarto errato. Training e test dovrebbero
essere separati anche per giorno di calibrazione o backend, non solo per
seed, e la soglia decisionale scelta su validation. Prima di addestrare
si deve infine verificare il bilanciamento della classe e confrontare il
modello con regole semplici come massa sui picchi, entropia e TOP-K.

Il contributo più credibile dell'apprendimento resta comunque a monte
della misura finale: sulla sindrome, dove può sfruttare correlazioni che
un decoder analitico non modella, anziché tentare di recuperare
informazione già persa nell'output.

\end{onehalfspace}
